\subsection{Two exact tests of the remaining nonflat coefficient route}

This continuation of \fid{F59} tests the actual finite
coinvariant quotient and the singular diagonal-invariant
correspondence. The first has a bounded bifree quadratic
with coefficient Euler class \(\delta^2\), nonzero total
Euler dimension, and explicitly nonprojective cohomology.
Its generator class is demonstrably different from the
designated nonstandard class. The second has nonzero Tor
in every positive degree, so its full derived fiber product
cannot be represented by a bounded complex.
Neither statement identifies global and generic top
cohomological degrees in the remaining nonflat case.

\subsubsection{A finite Frobenius countermodel with the wrong generator class}

Let \(k\) have characteristic zero, and put
\[
 A=k[x,y],\qquad \deg x=\deg y=2,\qquad
 S_0=k[x^2,y^2],\qquad B=A/(x^2,y^2).
\]
The two simple reflections negate \(x\) and \(y\)
separately. Thus \(B\) is their actual central
coinvariant quotient, with basis \(1,x,y,xy\) and
\[
 h_B(v)=\operatorname{Hilb}_v(B)=(1+v^2)^2.
\]
The coefficient of \(xy\) gives a nondegenerate
symmetric Frobenius pairing, homogeneous of degree
minus four: it pairs \(1\) with \(xy\), and \(x\)
with \(y\). Write \(M\langle e\rangle\) for a shift
of displayed internal degrees upward by \(e\).
For finite graded modules this pairing gives
\begin{equation}\label{eq:QRBdual}
 B\simeq\operatorname{Hom}_k(B,k)\langle4\rangle,
 \qquad
 \operatorname{Hom}_B(M,B)
       \simeq\operatorname{Hom}_k(M,k)\langle4\rangle.
\end{equation}
In particular, the \(B\)-dual is exact.

The ordinary Soergel convention is
\(M(1)^i=M^{i+1}\), so \((1)=\langle-1\rangle\).
The normalized reflection object is
\(B_s=A\otimes_{A^s}A(1)\), with cyclic generator
of degree minus one
\cite[\S3.3, p.~1105; \S3.4, p.~1106]{EW}.
For the two independent reflections,
\[
 B_1\otimes_A B_2\simeq(A\otimes_{S_0}A)(2).
\]
One can also see this directly by eliminating the
middle coordinates: the first reflection changes
only \(x\), and the second only \(y\).
The normalized longest object in the central
quotient is therefore
\begin{equation}\label{eq:QRlongest}
                  L_0=(B\otimes_k B)\langle-2\rangle.
\end{equation}
The shift below is this actual quotient normalization.

Define a finite free left \(B\)-complex by
\[
 K=[B\langle2\rangle\xrightarrow{x}B]
       \quad\text{in cohomological degrees }-1,0,
 \qquad L=B\oplus K\oplus K[1].
\]
Let \(L^\vee=\operatorname{Hom}_B(L,B)\), with the
usual cochain-dual differential signs. This is a
finite free right \(B\)-complex. Set
\begin{equation}\label{eq:QRZE}
 Z=(L\otimes_k L^\vee)\langle-2\rangle,
 \qquad
 E=(L^\vee\otimes_B L)\langle-2\rangle.
\end{equation}
Here \(E\) is a bounded complex of finite graded
\(k\)-vector spaces, used as a scalar coefficient.
Each term of \(Z\) is finite free on both sides.

\begin{finding}[F59: finite-quotient numerical rigidity fails]\label{prop:QRfinitecountermodel}
There is an actual chain isomorphism, hence a
derived isomorphism,
\begin{equation}\label{eq:QRquadratic}
        Z\otimes_B^{\mathbf L}Z\simeq E\otimes_k Z.
\end{equation}
Its Euler data are
\begin{equation}\label{eq:QREuler}
 \chi_v(E)=v^{-2}h_B(v)=v^2+2+v^{-2}=\delta^2,
 \qquad \chi_k(Z)=16=4\dim_k B.
\end{equation}
Nevertheless \(Z\) has nonprojective cohomology
as a left module and as a right module.
\end{finding}

\begin{proof}
All required derived tensors are computed by the
displayed bounded free complexes. Reassociation gives
\[
\begin{split}
 Z\otimes_B Z
 &=(L\otimes_k L^\vee\otimes_B L\otimes_k L^\vee)
       \langle-4\rangle\\
 &\simeq (L^\vee\otimes_B L)\langle-2\rangle
       \otimes_k(L\otimes_k L^\vee)\langle-2\rangle.
\end{split}
\]
The second comparison moves the middle finite
\(k\)-complex past the first factor with the usual
cohomological Koszul sign. It is a chain map and
is \(B\)-bimodule-linear: the outer actions remain
on \(L\) and \(L^\vee\). Its inverse is the reverse
symmetry/reassociation. This proves
\eqref{eq:QRquadratic} as an actual comparison.

The termwise graded free-rank Euler polynomial of
\(L\) is one, because \(K\) and \(K[1]\) cancel.
The same is true of its \(B\)-dual, after internal
degrees are inverted. Consequently the finite
\(k\)-Euler polynomial of \(L^\vee\otimes_B L\)
is \(h_B(v)\), and that of \(L\otimes_k L^\vee\)
is \(h_B(v)^2\). The shift in \eqref{eq:QRZE}
proves \eqref{eq:QREuler}. No acyclicity is inferred
from these Euler cancellations.

For explicit cohomology, put \(T=B/(x)\).
Since \(\operatorname{ann}_B(x)=(x)\), and the
source generator in \(K\) has degree two,
\[
 H^{-1}(K)=T\langle4\rangle,\qquad H^0(K)=T.
\]
It follows that
\[
\begin{array}{c|ccc}
 i&-2&-1&0\\\hline
 H^i(L)&T\langle4\rangle&
 T\langle4\rangle\oplus T&B\oplus T.
\end{array}
\]
Exact duality \eqref{eq:QRBdual}, together with
\(\operatorname{Hom}_B(T,B)=(x)=T\langle2\rangle\),
gives
\[
\begin{array}{c|ccc}
 i&0&1&2\\\hline
 H^i(L^\vee)&B\oplus T\langle2\rangle&
 T\langle-2\rangle\oplus T\langle2\rangle&
 T\langle-2\rangle.
\end{array}
\]
K\"unneth over the field \(k\) is exact.
Only the pair \((-2,0)\) contributes to
\(H^{-2}(Z)\), so
\[
 H^{-2}(Z)=
 \bigl(T\langle4\rangle\otimes_k
       (B\oplus T\langle2\rangle)\bigr)\langle-2\rangle.
\]
As a left \(B\)-module, this is
\begin{equation}\label{eq:QRleftnonflat}
 T\langle2\rangle\oplus
 T\langle4\rangle^{\oplus3}\oplus
 T\langle6\rangle^{\oplus2}.
\end{equation}
It is nonzero of dimension twelve and is annihilated
by left \(x\). A finite projective module over the
local algebra \(B\) is free: lift a residue basis,
split the resulting free surjection, and use
Nakayama on the complementary summand.
Nonzero \(x\) acts nontrivially on every nonzero
free \(B\)-module. Thus \eqref{eq:QRleftnonflat}
is not projective, and hence not flat.
Similarly, as a right module the degree-two group is
\begin{equation}\label{eq:QRrightnonflat}
 H^2(Z)=T\langle-4\rangle^{\oplus2}\oplus
        T\langle-2\rangle^{\oplus3}\oplus T,
\end{equation}
which is nonzero and annihilated by right \(x\).
It is nonprojective on that side.
The ungraded cohomology dimensions of \(Z\) are
\((12,32,56,32,12)\) in degrees \(-2,-1,0,1,2\):
they are the convolution of \((2,4,6)\) and
\((6,4,2)\). Their alternating sum is sixteen,
as required.
\end{proof}

\paragraph{An additive invariant separates the actual reduced target.}
Every term of \(Z\) is a sum of shifts of
\(B\otimes_k B\). The two free-rank Euler
polynomials used above are one, so its termwise
split Grothendieck Euler class is \([L_0]\).
The actual reduced target is
\begin{equation}\label{eq:QRtarget}
 \overline P=2[L_0]-\delta[\overline B_1]
                    -\delta[\overline B_2]+\delta^2[B],
\end{equation}
where \(\overline B_i\) denotes the central quotient
of the ordinary reflection bimodule. The classes
remain different after reduction; no injectivity
of a reduction map on Grothendieck groups is needed.

Indeed, use the additive graded invariant
\[
 \mathcal H(M)=\operatorname{Hilb}_v
                  \operatorname{Hom}_{B\text{-}B}^{\bullet}(M,B).
\]
For each of the unit, one-reflection and longest
cyclic bimodules, any image in the commutative ring
\(B\) defines a bimodule map from its cyclic generator.
The defining relations hold because the two actions
on the target \(B\) are equal. The generator degrees
therefore give
\[
 \mathcal H(B)=h_B,
 \qquad\mathcal H(\overline B_i)=v h_B,
 \qquad\mathcal H(L_0)=v^2h_B.
\]
A source shift \((1)\) raises the degree of such a
map by one, so multiplication by the source Laurent
variable is multiplication by \(v\) for this invariant.
It follows that
\begin{equation}\label{eq:QRclassseparation}
 \mathcal H(\overline P-[L_0])
    =(v^2-2\delta v+\delta^2)h_B
    =v^{-2}h_B=\delta^2\ne0.
\end{equation}
Thus Finding~\ref{prop:QRfinitecountermodel}
satisfies bounded bifreeness, the scalar-complex
quadratic, coefficient Euler class \(\delta^2\),
and nonzero total Euler dimension \(4\dim_k B\),
while its cohomology is not projective.
Those numerical hypotheses alone cannot establish
the desired rigidity. A result for the designated
target must retain additional class or action data.

\subsubsection{Every positive Tor group of the singular correspondence}

Keep \(A=k[x,y]\), but use the different invariant ring
\begin{equation}\label{eq:QRsingularring}
 S_\tau=A^\tau=k[a,b,c]/(ac-b^2),
 \qquad a=x^2,\ b=xy,\ c=y^2,
 \qquad \tau(x,y)=(-x,-y).
\end{equation}
The generators \(a,b,c\) have internal degree four.
Reducing the exponent of \(b\) to zero or one
gives distinct monomials in \(x,y\), precisely the
monomials of even total polynomial degree; this
proves the displayed invariant-ring presentation.
As an \(S_\tau\)-module, write
\(A=S_\tau\oplus M\) for its even and odd parts.
The odd part is generated by \(x,y\), of degree two,
with the complete relation matrix
\begin{equation}\label{eq:QRmatrix}
            \Phi=\begin{pmatrix}b&c\\-a&-b\end{pmatrix}.
\end{equation}
To check completeness, a relation \(px+qy=0\)
with even \(p,q\) has \(p=yh,q=-xh\) by coprimality
in \(k[x,y]\). The polynomial \(h\) is odd, so
\(h=xu+yv\) with even \(u,v\); these are exactly
the two columns of \eqref{eq:QRmatrix}.

\begin{finding}[F59: the full singular correspondence is unbounded]\label{prop:QRallTor}
For every \(i\geq1\),
\begin{equation}\label{eq:QRallTor}
             \operatorname{Tor}^{S_\tau}_i(A,A)
                    \simeq k\langle4+4i\rangle.
\end{equation}
Consequently \(A\otimes^{\mathbf L}_{S_\tau}A\)
has nonzero cohomology in every negative degree
and has no bounded derived replacement.
\end{finding}

\begin{proof}
Over the polynomial ring \(R=k[a,b,c]\), put
\(f=ac-b^2\). Direct multiplication gives
\[
                         \Phi^2=-fI.
\]
Thus reduction modulo \(f\) gives a complex with
\(\Phi\) at every positive differential. It is exact:
if \(w\in S_\tau^2\) lies in the kernel, lift it
to \(\widetilde w\in R^2\), and write
\(\Phi\widetilde w=fz\). Multiplying again gives
\(-f\widetilde w=f\Phi z\). The polynomial ring is
a domain, so cancellation yields
\(\widetilde w=-\Phi z\). Reduction puts \(w\)
in the image. The reverse containment follows
from \(\Phi^2=-fI\).
Together with the proved presentation of \(M\),
this gives the free resolution
\[
 \cdots\xrightarrow{\Phi}S_\tau\langle10\rangle^2
 \xrightarrow{\Phi}S_\tau\langle6\rangle^2
 \xrightarrow{\Phi}S_\tau\langle2\rangle^2
 \longrightarrow M\longrightarrow0.
\]
Its \(i\)-th free term has shift \(2+4i\).

After tensoring with \(A\), the matrix becomes
\[
 \Phi_A=
 \begin{pmatrix}xy&y^2\\-x^2&-xy\end{pmatrix}
      =\begin{pmatrix}y\\-x\end{pmatrix}(x\ \ y).
\]
Over the domain \(A\), its kernel is
\(A(y,-x)^t\), by coprimality applied to
\(xu+yv=0\). Its image is
\((x,y)A(y,-x)^t\). Every positive homology group
of this tensor resolution is therefore one copy
of \(A/(x,y)=k\).
In the \(i\)-th term, its kernel generator has
the two further internal degrees contributed
by \(x\) or \(y\), giving \(4+4i\).
Finally \(A=S_\tau\oplus M\), and the first
summand has no higher Tor. This proves
\eqref{eq:QRallTor} for \(\operatorname{Tor}(A,A)\).
Derived tensor cohomology in degree \(-i\) is
exactly this nonzero Tor group.
\end{proof}

With fixed spectator polynomial variables, the
same proof gives the shifted module \(A/(x,y)\)
in every positive Tor degree.
The singular ring \(S_\tau\) in
\eqref{eq:QRsingularring} is not the regular
full-product invariant ring \(S_0=k[x^2,y^2]\).
The ordinary module \(A\otimes_{S_\tau}A\), with
its previously computed torsion, is also a
different object from this full derived tensor.
Replacing it by the derived fiber product
introduces infinitely many nonzero groups;
reassociation does not remove them.
No arithmetic purity statement follows merely
from their displayed internal grading.

\paragraph{What these two tests leave open.}
The finite-Frobenius countermodel rules out
rigidity based only on the stated numerical
Euler data; \eqref{eq:QRclassseparation} explains
why it does not realize the designated target.
The Tor calculation excludes a bounded realization
of the particular full derived singular
correspondence. It does not exclude every other
bounded nonflat/nonflat ordinary Soergel
construction with a central polynomial coefficient.
The remaining target still requires the designated
ambient Euler class, actual coherent central
movement when imposed, and an actual bounded
quadratic. No equality of generic and global
top cohomological degrees is asserted for it.

% Separate paragraph for Part III.
\paragraph{The F59 quotient tests in Grothendieck's language.}
The coinvariant quotient, cyclic bimodules, and
singular invariant correspondence are actual
affine module data. EGA~I, Theorem~1.3.7, p.~86,
Corollary~1.3.8, p.~87, and Theorem~1.4.1,
pp.~90--92 give their coherent-sheaf dictionary
\cite{EGA}. The explicit Tor groups distinguish
an underived fiber product from its derived
version; the explicit finite-Frobenius example
distinguishes nonzero Euler dimension from
projective cohomology.
SGA~6, I, Definition~4.7, pp.~121--122, and
Example~4.8, p.~122 provide the bounded perfect
complex framework \cite{SGA6}. The complex \(Z\)
is perfect on each side, whereas the full derived
object in \eqref{eq:QRallTor} is unbounded.
Its scalar coefficient \(E\) is a finite complex
of actual vector spaces, not a rational virtual
replacement. Internal shifts use SGA~3, I,
4.7.2--3/cote~90, scan~17; split and virtual
classes use SGA~6, V, \S\S2--3/cote~133, scan~18
\cite{SGA3,SGA6,Archive}, with the existing
source-access qualification.
No noncommutative \(\lambda\)-ring, arithmetic
purity, or designated canonical-label comparison
is inferred from these calculations.
